\documentclass[12pt]{article}


% --- math/fonts ---
\usepackage{amsmath,amssymb}
\usepackage{braket}
\usepackage[T1]{fontenc}
\usepackage{libertine}
\usepackage[utf8]{inputenc}

% --- graphics/tikz ---
\usepackage{tikz}
\usepackage{qcircuit}

% --- other stuff ---
\usepackage[margin=1in]{geometry}
\usepackage{pdfpages}
\usepackage{multicol}
\usepackage[shortlabels]{enumitem}
\usepackage{cancel}
\usepackage{graphicx}
\usepackage{float}
\usepackage{listings}
\lstset{
    language=Python,
    basicstyle=\ttfamily\footnotesize,
    keywordstyle=\bfseries,
    frame=single,
    breaklines=true,
    keepspaces=true,
    columns=fullflexible,
    showstringspaces=false
}
\let\epsilon\varepsilon
\usepackage{enumitem}
\setlist[enumerate]{label=\arabic*.}
\usepackage[authoryear,round]{natbib}
\usepackage[hidelinks]{hyperref}


\newcommand{\class}{CSCE883 Machine Learning}           
\newcommand{\examnum}{Homework 2}      
\newcommand{\examdate}{09/26/2026}        
\newcommand{\timelimit}{}             




\begin{document} 
\pagestyle{plain}
\thispagestyle{empty}

\noindent
\begin{tabular*}{\textwidth}{l @{\extracolsep{\fill}} r @{\extracolsep{6pt}} l}
 \textbf{\class} & \textbf{Name:} & \textit{Taylor Belcher}\\       
\textbf{\examnum} &&\\
\textbf{\examdate} &&\\
\end{tabular*}\\
\rule[2ex]{\textwidth}{2pt}
% ---

 
\begin{enumerate}
   \item \emph{(Question 1).} Let $X$ be a discrete random variable with the given probability mass function with $0 \leq \theta \leq 1$ as a parameter:
   
  \[
\begin{array}{|c|c|c|c|c|}
\hline
x & 0 & 1 & 2 & 3 \\
\hline
P(X=x\mid\theta) & 2\theta/3 & \theta/3 & 2(1-\theta)/3 & (1-\theta)/3 \\
\hline
\end{array}
\] 

We have the following i.i.d. sample $\mathbf{x} = (3,0,2,1,3,2,1,0,2,1)$ from $X$ and we want the maximum likelihood estimate of $\theta$. Let $x^t$ denote the $t$-th observation, for $t=1,\ldots,N$, where here $N=10$, so that $\mathbf{x}=(x^1,\ldots,x^N)$. Since the observations are independent, the likelihood of the parameter $\theta$ is 

\[
\begin{aligned}
l(\theta \mid \mathbf{x}) &= \prod\limits_{t=1}^{N} p(x^{t} \mid \theta) \\
&= p(3 \mid \theta) \cdot p(0 \mid \theta) \cdot p(2 \mid \theta) \cdot p(1 \mid \theta) \cdot p(3 \mid \theta) \\
&\quad {}\cdot p(2 \mid \theta) \cdot p(1 \mid \theta) \cdot p(0 \mid \theta) \cdot p(2 \mid \theta) \cdot p(1 \mid \theta) \\
&= \frac{1-\theta}{3} \cdot \frac{2\theta}{3} \cdot \frac{2(1-\theta)}{3} \cdot \frac{\theta}{3} \cdot \frac{1-\theta}{3} \\
&\quad {}\cdot \frac{2(1-\theta)}{3} \cdot \frac{\theta}{3} \cdot \frac{2\theta}{3} \cdot \frac{2(1-\theta)}{3} \cdot \frac{\theta}{3} \\
&= \left(\frac{2\theta}{3}\right)^2
   \left(\frac{\theta}{3}\right)^3
   \left(\frac{2(1-\theta)}{3}\right)^3
   \left(\frac{1-\theta}{3}\right)^2 \\
&= \frac{2^5}{3^{10}}\theta^5(1-\theta)^5.
\end{aligned}
\] 

We can optimize via the log likelihood and our desired MLE for $\theta$ is $\hat{\theta} = \operatorname*{arg\,max}\limits_{\theta}\left( L(\theta \mid \mathbf{x})\right)$ so consider 

\[
\begin{aligned}
L(\theta \mid \mathbf{x})
&= \log\left(l(\theta \mid \mathbf{x})\right) \\
&= \log\left(\frac{2^5}{3^{10}}\theta^5(1-\theta)^5\right).
\end{aligned}
\]

Applying rules of logarithms we have

\[
L(\theta \mid \mathbf{x}) = \log\left(\frac{2^5}{3^{10}}\right) + 5 \log \theta + 5 \log(1-\theta).
\]

Taking the derivative of $L$ with respect to $\theta$

\[
\begin{aligned}
\frac{dL}{d\theta}
&= 0 + 5\cdot \frac{1}{\theta} + 5\cdot \frac{1}{1-\theta}\cdot (-1) \\
&= \frac{5}{\theta} - \frac{5}{1-\theta}.
\end{aligned}
\]

Set the derivative equal to zero and solve for $\theta$:

\[
\begin{aligned}
0 &= \frac{5}{\theta} - \frac{5}{1-\theta} \\
\frac{5}{1-\theta} &= \frac{5}{\theta} \\
\theta &= 1 - \theta \\
\theta &= \frac{1}{2}.
\end{aligned}
\]

To verify that this gives a maximum, rewrite the derivative with a common denominator:

\[
\begin{aligned}
\frac{dL}{d\theta}
&= \frac{5}{\theta} - \frac{5}{1-\theta} \\
&= \frac{5(1-\theta)-5\theta}{\theta(1-\theta)} \\
&= \frac{5(1-2\theta)}{\theta(1-\theta)}.
\end{aligned}
\]

For $0<\theta<1$, the denominator is positive, so the sign of the derivative is determined by $1-2\theta$. Thus the derivative is positive for $0<\theta<\frac{1}{2}$, zero at $\theta=\frac{1}{2}$, and negative for $\frac{1}{2}<\theta<1$. The log likelihood therefore increases up to $\theta=\frac{1}{2}$ and decreases afterward. Since the logarithm is strictly increasing, this also maximizes the likelihood on $(0,1)$. At the endpoints $\theta=0$ and $\theta=1$, the likelihood is zero, so neither endpoint gives a larger value.

Therefore the maximum likelihood estimate of $\theta$ is $\hat{\theta} = \boxed{\frac{1}{2}}$. 

\item \emph{(Question 2).} Let $X_1, X_2, \ldots, X_N$ be i.i.d. random variables with density function
\[
f(x \mid \sigma) = \frac{1}{2\sigma}e^{-\frac{|x|}{\sigma}}.
\]

Again, since the observations are independent, we define the likelihood function as

\[
\begin{aligned}
l(\sigma \mid x) \equiv p(x \mid \sigma)
&= \prod_{t=1}^N p(x_t \mid \sigma) \\
&= \prod_{t=1}^N f(x_t \mid \sigma) \\
&= \prod_{t=1}^N \left(\frac{1}{2\sigma}e^{-\frac{|x_t|}{\sigma}}\right) \\
&= \left(\frac{1}{2\sigma}\right)^N
   \prod_{t=1}^N \left(e^{-\frac{|x_t|}{\sigma}}\right) \\
&= \left(\frac{1}{2\sigma}\right)^N
   e^{\left(\displaystyle \sum_{t=1}^N -\frac{|x_t|}{\sigma}\right)}\\
&= \left(\frac{1}{2\sigma}\right)^N
   e^{\left(\displaystyle -\frac{1}{\sigma} \sum_{t=1}^N |x_t|\right)}.
\end{aligned}
\] 

We then take the log, which we will optimize.

\[
\begin{aligned}
L(\sigma \mid x)
&= \log\left(l(\sigma \mid x)\right) \\
&= \log\left[\left(\frac{1}{2\sigma}\right)^N
   e^{-\frac{1}{\sigma}\sum_{t=1}^N |x_t|}\right].
\end{aligned}
\]

Applying rules of logarithms we have

\[
\begin{aligned}
L(\sigma \mid x)
&= \log\left[\left(\frac{1}{2\sigma}\right)^N\right]
   + \log\left[e^{-\frac{1}{\sigma}\sum_{t=1}^N |x_t|}\right] \\
&= N\log\left(\frac{1}{2\sigma}\right)
   - \frac{1}{\sigma}\sum_{t=1}^N |x_t| \\
&= -N\log(2\sigma) - \frac{1}{\sigma}\sum_{t=1}^N |x_t| \\
&= -N\log 2 - N\log\sigma - \frac{1}{\sigma}\sum_{t=1}^N |x_t|.
\end{aligned}
\]

Taking the derivative of $L$ with respect to $\sigma$

\[
\begin{aligned}
\frac{dL}{d\sigma}
&= 0 - N\cdot\frac{1}{\sigma}
   - \left(-\frac{1}{\sigma^2}\right)\sum_{t=1}^N |x_t| \\
&= -\frac{N}{\sigma} + \frac{1}{\sigma^2}\sum_{t=1}^N |x_t|.
\end{aligned}
\]

Set the derivative equal to zero and solve for $\sigma$:

\begin{align*}
0 &= -\frac{N}{\sigma} + \frac{1}{\sigma^2}\sum_{t=1}^N |x_t| \\
\frac{N}{\sigma} &= \frac{1}{\sigma^2}\sum_{t=1}^N |x_t| \\
\intertext{Since $\sigma > 0$, we have $\sigma^2 > 0$, so we can multiply both sides by $\sigma^2$:}
N\sigma &= \sum_{t=1}^N |x_t| \\
\sigma &= \frac{1}{N}\sum_{t=1}^N |x_t|.
\end{align*}

To verify that this gives a maximum, let $\hat{\sigma}=\frac{1}{N}\sum_{t=1}^N |x_t|$ and rewrite the derivative with a common denominator:

\[
\begin{aligned}
\frac{dL}{d\sigma}
&= -\frac{N}{\sigma} + \frac{1}{\sigma^2}\sum_{t=1}^N |x_t| \\
&= \frac{\sum_{t=1}^N |x_t|-N\sigma}{\sigma^2} \\
&= \frac{N\hat{\sigma}-N\sigma}{\sigma^2} \\
&= \frac{N(\hat{\sigma}-\sigma)}{\sigma^2}.
\end{aligned}
\]

For $\sigma>0$, the factor $N/\sigma^2$ is positive, so the sign of the derivative is determined by $\hat{\sigma}-\sigma$. When $\sum_{t=1}^N |x_t|>0$, we have $\hat{\sigma}>0$. Thus the derivative is positive for $0<\sigma<\hat{\sigma}$, zero at $\sigma=\hat{\sigma}$, and negative for $\sigma>\hat{\sigma}$. The log likelihood therefore increases up to $\hat{\sigma}$ and decreases afterward, so this gives the maximum. Since the logarithm is strictly increasing, maximizing the likelihood is equivalent to maximizing its log. Therefore the maximum likelihood estimate of $\sigma$ is

\[
\begin{aligned}
\hat{\sigma}
&= \operatorname*{arg\,max}_{\sigma > 0} l(x \mid \sigma) \\
&= \operatorname*{arg\,max}_{\sigma > 0} \log\left(l(x \mid \sigma)\right) \\
&= \operatorname*{arg\,max}_{\sigma > 0} L(\sigma \mid x) \\
&= \boxed{\frac{1}{N}\sum_{t=1}^N |x_t|}.
\end{aligned}
\]


\end{enumerate}

\end{document}

% Quick aligned table for equations and explanations
%
%\begin{center}
%  \begin{tabular}{l l l l}
%    % template line 
%    % LHS & $=$ & RHS & \emph{(comment)}\\
%   \end{tabular}
% \end{center}
